\section{Why The Mathematical Layer Must Grow, But Not Drift}

The early versions of this manuscript used a potential-well picture because it
was teachable. That was appropriate for a first pass, but it is no longer
enough. Once the book begins to talk seriously about conscious-state
transitions, framework formation, and intervention design, the mathematical
layer must become more explicit about what kind of system is being imagined.

At the same time, this chapter must stay under evidence discipline. Most of the
content here belongs to the \textbf{medium} lane. The chapter uses the
literature map's advanced-dynamics sources to improve explanatory precision:
the Nature Reviews Neuroscience review on attractor and integrator networks in
the brain, the metastability review already used in Chapter~\ref{ch:emergence},
and the local indexed arXiv papers on self-organizing attractor networks and
Koopman-style viewpoints. These sources help the manuscript speak more clearly
about persistence, switching, oscillation, and feedback. They do \emph{not}
prove that a human daily routine is literally a low-dimensional differential
equation with known parameters.

\begin{discussion}
The business value of mathematical depth is therefore narrow and concrete. It
should answer questions such as:
\begin{enumerate}[nosep]
  \item what makes a state locally stable;
  \item what makes a transition abrupt rather than gradual;
  \item why some patterns recur rhythmically;
  \item why repeated history can make later re-entry easier;
  \item how interventions can be interpreted as control inputs rather than
        moral exhortations.
\end{enumerate}
If a mathematical tool does not sharpen one of those questions, it does not
belong in the main chapter.
\end{discussion}
